% Proposed B.126. Isolated owned fragment; central assembly unchanged.
\subsection*{The exact face and axes of the final prefix relations}
\statement{B.126}{Proved}{The intermediate corner has the exact completed face, while the four final roots form a discrete malnormal cyclic family with sharp two-edge axis overlaps}
\label{stmt:B126}
Use the actual intermediate amalgam of B.120:
\[
 M=R*_{H}W_\partial,
 \qquad H=\langle c,k,T\rangle=C_2(c)*\mathbb Z^2(k,T),
 \qquad P=P_5=cT.
\]
The corner satisfies $j^2=1$, $jPj=P^{-1}$ and
$[jcj,k]=1$. Its mirror-even subgroup has the exact path
presentation on $D,k,u,k',c$, with $D=jcj$, $k'=jkj$,
$T=uD$, and edges $D$--$k$--$u$--$k'$--$c$.
Then
\begin{equation}\label{eq:B126-corner-face}
 H\cap jHj=\langle P\rangle,
 \qquad
 \widehat H\cap j\widehat H j
      =\overline{\langle P\rangle}
       \quad\hbox{inside }\widehat W_\partial.
\end{equation}
The completed subgroups here are the actual full embedded completions
from B.118--B.120.

Let $\mathcal T_M$ be the ordinary Bass--Serre tree, let $e=H$
be its base edge, and for $1\leq i\leq4$ put
\[
 w_i=jP_i,\qquad \mathcal A_i=\operatorname{Axis}_{\mathcal T_M}(w_i).
\]
Each $w_i$ has translation length $2$. Distinct translates of
these four axes intersect in at most two edges, and that bound
is attained. Moreover
\begin{equation}\label{eq:B126-axis-stabilizers}
 \operatorname{Stab}_M(\mathcal A_i)
  =C_M(w_i)=N_M(\langle w_i\rangle)=\langle w_i\rangle.
\end{equation}
The four cyclic groups form a malnormal family: if
\[
 \langle w_i\rangle\cap g\langle w_l\rangle g^{-1}\ne1,
\]
then $i=l$ and $g\in\langle w_i\rangle$.
These axis statements concern the ordinary tree and discrete $M$.
They are not asserted for the full completed action or after the
four actual relations $w_i^2=1$ are imposed.

\begin{proof}
\emph{B.126.1. The actual corner face in both categories.}
The mirror-even corner has the retraction of B.120
\[
 \rho:W_\partial^0\longrightarrow H,
 \qquad \rho(D)=\rho(k')=1,\quad\rho(u)=T,
 \quad\rho(c)=c,\quad\rho(k)=k.
\]
Its restriction to $jHj$, after identifying that subgroup with $H$,
is the companion map
\[
 \nu(c)=\nu(k)=1,\quad\nu(T)=P^{-1},\quad\nu(P)=P^{-1}.
\]
Thus $\rho(jHj)=\langle P\rangle$. Since $\rho$ fixes $H$,
every element of $H\cap jHj$ lies in $\langle P\rangle$.
The relation $jPj=P^{-1}$ gives the reverse inclusion.

The same argument is continuous. The mirror-even corner is open;
its completion is the actual open subgroup of
$\widehat W_\partial$. Its completed retraction fixes
$\widehat H$ and maps $j\widehat H j$ onto the actual procyclic
closure of $P$. This proves the second equality in
\eqref{eq:B126-corner-face} directly. It does not exchange
completion with an arbitrary subgroup intersection.

\emph{B.126.2. The common turn and its full ordinary stabilizer.}
By B.118, $K\cap H=\langle P\rangle$, and
$P_1,\ldots,P_5$ are the literal free basis of $K$.
Hence $P_i\notin H$ for $i\leq4$. Also $j\notin H$ by mirror
parity. Thus $w_i=jP_i$ is cyclically reduced of amalgam syllable
length two, with ordinary translation length two.
Its axis contains the four consecutive edges
\begin{equation}\label{eq:B126-four-edges}
 P_i^{-1}e,\quad e,\quad je,\quad jP_i e.
\end{equation}
Their five vertices are
$P_i^{-1}W_\partial,R,W_\partial,jR,jP_iW_\partial$.
The element $w_i$ sends the first two edges to the last two;
ordinary amalgam normal form makes all successive repetitions
geodesic. In particular every $W_\partial$-type turn of the axis
is a power-of-$w_i$ translate of $e\cup je$.

The pointwise stabilizer of this turn is exactly
$H\cap jHj=\langle P\rangle$. An element preserving the unordered
turn fixes its middle $W_\partial$ vertex. If it switches the two
edges, multiplication by $j$ reduces to the pointwise case.
Therefore the setwise stabilizer is precisely
\begin{equation}\label{eq:B126-turn-stabilizer}
 Q_0=\langle P,j\rangle
     =\langle P\rangle\sqcup j\langle P\rangle.
\end{equation}

\emph{B.126.3. All coefficient cases for a three-edge overlap.}
Suppose $g\mathcal A_i$ and $\mathcal A_l$ have at least three
common edges. Their common path contains a turn centered at a
$W_\partial$ vertex. Multiplying $g$ on the left and right by
appropriate powers of $w_l,w_i$ aligns both turns with
$e\cup je$, without changing the axes. The resulting element
$a$ is either $P^n$ or $jP^n$, for an arbitrary integer $n$.
At least one outside edge of the turn also agrees. Equality of
edges is equality of left $H$ cosets, giving exactly these four
cases:
\begin{center}
\begin{tabular}{lll}
\toprule
$a$ & Extension & Required membership in $H$\\
\midrule
$P^n$ & At $R$ & $P_lP^nP_i^{-1}$\\
$P^n$ & At $jR$ & $P_l^{-1}P^{-n}P_i$\\
$jP^n$ & At $R$ & $P_lP^{-n}P_i$\\
$jP^n$ & At $jR$ & $P_l^{-1}P^nP_i^{-1}$\\
\bottomrule
\end{tabular}
\end{center}
For example, in the third row the image of the right outside edge
is $jP^n jP_i e=P^{-n}P_i e$, which must equal $P_l^{-1}e$.
For the second and fourth rows cancel the common initial $j$.
Only $jPj=P^{-1}$ and the explicit four edges
\eqref{eq:B126-four-edges} are used here.

Every displayed coefficient lies in $K$. The exact regional face
$K\cap H=\langle P\rangle$ therefore replaces membership in
$H$ by membership in $\langle P\rangle$. Now use free reduced
words in the actual basis $P_1,\ldots,P_5$. For every integer
$n$, the first two rows hold exactly when $i=l$ and $n=0$.
The last two never hold: the two non-$P$ letters have the same
sign, and even $n=0,i=l$ leaves $P_i^2$ or $P_i^{-2}$.
Thus $a=1$ and $i=l$. Undoing the two normalizing powers shows
the axes were the same. This proves the asserted upper bound.
For $i\ne l$, the axes share $e\cup je$ and neither outside
edge, so the bound is sharp.

\emph{B.126.4. Stabilizers, centralizers and malnormality.}
Apply the same normalization to an element preserving
$\mathcal A_i$. The three-edge comparison forces its normalized
value to be $1$; undoing the powers gives an element of
$\langle w_i\rangle$. This proves the first equality in
\eqref{eq:B126-axis-stabilizers}. In particular the pointwise
axis stabilizer is trivial and no element reverses the axis.
A centralizing or normalizing element preserves the unique axis
of the nontrivial hyperbolic element, which proves the other two
equalities.

If two conjugated cyclic groups have a nontrivial common element,
nonzero powers of their generators have the same axis. The preceding
three-edge calculation then gives $i=l$ and
$g\in\langle w_i\rangle$. This proves the complete discrete
malnormal-family statement.

\emph{B.126.5. Exact filling limitation.}
The actual remaining relator $w_i^2$ has translation length four,
whereas the sharp distinct-axis overlap bound is two. In particular
this is not a strict one-sixth or one-fourth metric small-cancellation
estimate. No polygon-link condition, developability theorem,
post-filling centralizer formula or profinite axis analogue follows
from this proof. The actual final neck embedding remains the
independent retraction of B.121, not a consequence of the overlap
calculation. All four final relations remain part of the main
fixed-source problem.
\end{proof}

\paragraph{Proof provenance.}
The coefficient calculation is uniform over every translate and every
integer power; no finite orbit or word-length cutoff is used.
\begingroup\raggedright
\path{work_A/final_prefix_axis_overlaps25.md};
\path{work_CE/review_final_prefix_axis_overlaps25.txt}.
\par\endgroup
